% New English translation from checked cote 25, site 22 (PDF page 23).
% AI-translated by GPT-6 Astra (OpenAI, ChatGPT Pro mode).
\scanpage{25/22}
\mspage{\unsure{\hand{19}|19}}

5. \uline{Theorems of Seidenberg type}\var{\xout{\unsure{z}}}{}{typewriter strike-out}{ega25_25_022-v0001}

5.1.
\var{In the present no., we give conditions .... (cf the preceding
page, the text struck out by mistake).}{\restoredfrom{25/21}{https://grothendieck.umontpellier.fr/web/viewer.html?file=../25.pdf\#page=22}{%
In \illegible{1} \hand{the} present n$^o$, we give conditions \unsure{under} which
the set E considered in 4.1. is open. \var{\xout{Pxxxxxxx}}{}{typewriter strike-out}{ega25_25_021-v0009} We
shall deal here with properties P of a local nature on X resp. $Y_{\hand{\xi}}$, so
that we can define the set U of $y \hand{\in} \underline{Y}$ such that (if \hand{$\xi$} denotes the image
of y in $\hand{\check{}}P$) $Y_{\hand{\xi}}$ satisfies P at the point y (resp. $G_{\hand{\xi}}$ satisfies
P at y). We shall first give criteria for U to be open,
using par.12. \var{\xout{xxxxxxxxxxxxxxxxxxxx}}{}{typewriter strike-out}{ega25_25_021-v0015} Since in any case we have
\[ E = P - h(\underline{Y} - U) , \]
it follows that if U is open and X proper over S (and hence h proper and
a fortiori closed \var{\inserted{since $\underline{Y} \to \unsure{Y}$ proper over X and $\overset{\hand{\vee}}{P}$ separated over S}}{since $\underline{Y} \to \unsure{Y}$ proper over X and $\overset{\hand{\vee}}{P}$ separated over S}{insertion}{ega25_25_021-v0020}), then E is also open.
}}{author-directed restoration}{ega25_25_022-v0012}

5.2. As we saw in no. 1, \underline{Y} is defined in $X x_S \overset{\hand{\vee}}{P} = X_{\overset{\hand{\vee}}{P}}$
as the “divisor” of a section $\phi$ of $\underline{O}_X(1)\otimes_S \underline{O}_{\overset{\hand{\vee}}{P}}(1)$, which induces for
every $\hand{\xi} \hand{\in} \overset{\hand{\vee}}{P}$ a section $\phi_{\hand{\xi}}$ of \var{\xout{$\otimes$xxxxx(1)$\otimes$xxxx xxxxxxxxxxxxxx k( )}}{}{typewriter strike-out}{ega25_25_022-v0002} $\underline{O}_X(1)\otimes_{k(s)}\underline{O}_{\overset{\hand{\vee}}{P}}(1)(\hand{\xi})$
(a sheaf moreover noncanonically isomorphic to $\underline{O}_X(1)\otimes_{k(s)}k(\hand{\xi}) =$
$\underline{O}_{X_{k(\hand{\xi})}}(1)$), so that $Y_{\hand{\xi}}$ is precisely the \var{\xout{xxxxx}}{}{typewriter strike-out}{ega25_25_022-v0003} “divisor” of
this section (NB the divisor of a section $\phi$ of an invertible module \underline{L} \var{\margin{$/\!/$}}{\placedmargin{$/\!/$}}{marginal note placed}{ega25_25_022-v0004}
is defined as the closed subprescheme defined by the Ideal that is the image of
$\phi^{-1} : \underline{L}^{-1} \hand{\longrightarrow} \underline{O}$). If \underline{F} is a Module on X, then its inverse image \underline{G}
on \underline{Y}, i.e. the inverse image of $\underline{F}\otimes_{\underline{O}_S}\underline{O}_{\overset{\hand{\vee}}{P}} = \underline{F}_{\overset{\hand{\vee}}{P}}$ on the subprescheme \underline{Y} of
$X_{\overset{\hand{\vee}}{P}}$, is precisely the cokernel of the homomorphism
\[ \phi^{-1}\otimes \mathrm{id}_{\underline{F}_{\overset{\hand{\vee}}{P}}} : \underline{F}_{\overset{\hand{\vee}}{P}}(-1,-1) \mathrel{\hand{-}} F_{\overset{\hand{\vee}}{P}} , \]
where the notation (-1,-1) explains itself, as Mike would say. Likewise,
$G_{\hand{\xi}}$ is the cokernel of the analogous homomorphism
\[ F_{k(\hand{\xi})}(-1\var{\over{)}{,}}{,}{writing over}{ega25_25_022-v0006}-1) \hand{\longrightarrow} F_{k(\hand{\xi})} \]
where $\hand{\xi}$ is a point of $\overset{\hand{\vee}}{P}$ (and of course we have a corresponding
interpretation when $\hand{\xi}$, instead of being a point of $\overset{\hand{\vee}}{P}$, denotes a point
of $\overset{\hand{\vee}}{P}$ with values in an S' over S \var{\xout{\unsure{??}}}{}{typewriter strike-out}{ega25_25_022-v0008}. ). \hand{\unsure{(}}In general, if L is an
invertible Module somewhere, $\phi$ a section defining the subprescheme
V($\phi$), then for every Module F, the inverse image of F on V($\phi$)
identifies, with the usual abuse of language, with the cokernel of $\mathrm{id}_F\otimes\phi$
$F\otimes L^{-1} \hand{\longrightarrow} F$; We say that $\phi$ is F-\uline{regular} if the preceding homomorphism
is injective. If we choose an isomorphism of F and $\underline{O}_X$, -- which
is possible locally -- so that $\phi$ identifies with a section of
$\underline{O}_X$, this terminology is compatible with that \var{\inserted{already}}{already}{insertion}{ega25_25_022-v0010} introduced elsewhere
\tnote{25-022-1}{The restored 5.1 is copied from the translated paragraph on site21, and its scan and doubt links point there. The abbreviated incipit and the instruction struck out by mistake remain in the apparatus. The source spellings sedtion, estb and conayau remain in the French and are rendered section, is and cokernel. The source names Mike without further identification. The letter F in the sentence choosing an isomorphism with the structure sheaf is retained, as are the hyphen in the first homomorphism display and all of its subscripts.}

\var{\margin{22}}{\placedmargin{22}}{marginal note placed}{ega25_25_022-v0011}
